% Debug chapter (D-series, 2026-08). CONFIGURATION-DRIVEN.
% The chapter itself always renders: what JTAG is and what the RISC-V debug
% stack does are architecture-level material that a reader of ANY build needs,
% and a build without the debug system still has something true to say -- that
% the four debug CSRs are illegal, that DRET is illegal, and that there is no
% debug port. That negative statement is the honest description of the default
% hardware, exactly as the Privileged Architecture chapter's legacy-trap
% opening is. The implementation sections are wrapped in \ifdebugenable,
% emitted by LatexUserGuide.GenerateDefinesFile() from debug.enable, so a build
% documents only the debug hardware it actually contains.
% Sources (do not write this chapter from memory of the RTL): the frozen
% D-series specifications d1_spec.md / d2_spec.md / d3_spec.md /
% d3_cdc_spec.md, the RTL headers of hdl/common/jtag_dtm.vhd and
% hdl/common/debug_module.vhd, and the phase reports.
\newpage

\section{Debug Support} \label{s:debug}

Every program eventually needs a way to ask the hardware what it is actually
doing. The chapters so far have described two: the Forth interpreter over
\peripheral{UART0} (Section \ref{s:forth}), which lets you poke at a running
system from a terminal, and the disassembly listings the compiler emits with
\texttt{-g}, which let you read what the machine was told to do. Both are
\textit{cooperative} --- they need the chip to still be executing your code,
correctly enough to talk back.

\textbf{Hardware debug is the non-cooperative kind.} It reaches into the chip
over dedicated wires that do not depend on any software running, stops a
processor between instructions, reads and writes its registers and memory, and
starts it again. It works on a hart that is stuck in an infinite loop, on a hart
that has taken an illegal instruction and stopped retiring, and --- with one
configuration bit set --- on a hart before it has executed its first
instruction. That is what this chapter is about.

\ifdebugenable
This configuration includes the debug system; the sections below document it.
\else
\textbf{This configuration does not include the debug system.} The four
Debug-Mode control registers \register{dcsr}, \register{dpc},
\register{dscratch0} and \register{dscratch1} (\texttt{0x7B0}--\texttt{0x7B3})
are not implemented: any CSR instruction naming an address in the range
\texttt{0x7B0}--\texttt{0x7BF} is an illegal instruction, from any privilege
level. The \asminline{DRET} instruction (\texttt{0x7B200073}) is an illegal
instruction. There is no JTAG port, no Debug Transport Module and no Debug
Module, and no external agent can halt, inspect or resume a hart. Sections
\ref{ss:debug-jtag} and \ref{ss:debug-stack} below describe the architecture
this family implements when it \textit{is} configured, and are included because
they are useful background; everything after them documents hardware this build
does not contain and is therefore omitted. Building a chip with
\texttt{debug.enable} set (Table \ref{t:chip-config}) adds all of it. Note that
the \texttt{-g} compiler and linker flag remains useful without it: it produces
the source-to-assembly correspondence in \texttt{.lst} listings, which is a
build-time facility and needs no hardware at all (Section \ref{s:software}).
\fi

\subsection{What JTAG Is} \label{ss:debug-jtag}

JTAG is a serial protocol standardised as IEEE 1149.1 --- four mandatory wires
plus an optional fifth --- that was invented to test the solder joints on a
circuit board and has since become the way almost every processor in the world
is debugged. Understanding it is worth twenty minutes, because everything above
it in the stack is built on one very simple idea.

The wires are:

\begin{center}
\begin{tabular}{ l l p{9.2cm} }
	\textbf{Wire} & \textbf{Direction} & \textbf{Purpose} \\ \hline
	\pin{TCK} & in & Test clock. Everything on the port is synchronous to it, and it is supplied by the debug probe, not by the chip. It has no relationship to any clock inside the chip. \\
	\rowcolor{tablehighlightcolor} \pin{TMS} & in & Test mode select. Sampled on every rising \pin{TCK} edge; this one wire drives the entire state machine. \\
	\pin{TDI} & in & Test data in --- the serial input to whichever shift register is currently selected. \\
	\rowcolor{tablehighlightcolor} \pin{TDO} & out & Test data out --- the serial output of that register. \\
	\pin{TRSTn} & in & The optional fifth wire: an asynchronous reset for the port. Not required by the standard, which can always reset the port with \pin{TMS} alone, but present here. \\ \hline
\end{tabular}
\end{center}

Inside the chip these wires drive a \textbf{Test Access Port}, or TAP: a
sixteen-state finite state machine whose only input is \pin{TMS}. The states
are fixed by the standard and every JTAG device in the world has the same
sixteen. They form two nearly identical lobes --- one for scanning an
\textit{instruction} and one for scanning \textit{data} --- joined at a common
idle state:

\begin{itemize}
	\item \textbf{Test-Logic-Reset} is the top of the graph and the reset state. \textbf{Holding \pin{TMS} high for five \pin{TCK} clocks reaches it from anywhere in the graph}, which is the standard's guaranteed recovery: a debugger that has lost track of where the state machine is does not need to know, it just clocks five ones.
	\item \textbf{Run-Test/Idle} is the resting state between operations.
	\item The \textbf{IR lobe} --- Select-IR, Capture-IR, Shift-IR, Exit1-IR, Pause-IR, Exit2-IR, Update-IR --- shifts a new value into the \textit{instruction register}.
	\item The \textbf{DR lobe} --- the same seven states with DR in place of IR --- shifts a value through whichever \textit{data register} the current instruction has selected.
\end{itemize}

Figure \ref{fig:tap-state-graph} is the whole machine, with every transition this chip implements.

% Generated by LatexUserGuide.GenerateTapStateDiagram (make chip).
\begin{figure}[htpb]
	\centering
	\resizebox{\linewidth}{!}{\input{include/TapStateDiagram.tex}}
	\caption[{The sixteen-state JTAG test access port.}]{The sixteen-state JTAG test access port. Solid arrows are taken when \pin{TMS} is sampled low, dashed when it is sampled high. The two lobes hanging below Select-DR and Select-IR are the same seven states drawn as mirror images of one another, one scanning a data register and one the instruction register. Note that five consecutive high samples reach Test-Logic-Reset from every state in the graph, which is the recovery a debugger uses when it does not know where the machine is.}
	\label{fig:tap-state-graph}
\end{figure}

The operation you actually perform is called a \textbf{scan}, and it has three
phases that matter. In \textbf{Capture}, the selected register is loaded in
parallel from whatever it is shadowing. In \textbf{Shift}, the register behaves
as a shift chain: on each \pin{TCK} it takes one bit from \pin{TDI} and presents
one bit on \pin{TDO}, least-significant bit first. In \textbf{Update}, the value
that has just been shifted in is committed. \textbf{So one scan simultaneously
reads the old contents and writes the new ones}, which is why a debugger's
traffic always looks like a stream of fixed-length shifts rather than the
read-then-write pairs a parallel bus would use. Figure \ref{fig:jtag-scan}
follows one four-bit scan through those phases.

% Generated by LatexUserGuide.GenerateJtagScanDiagram (make chip).
\begin{figure}[htpb]
	\centering
	\input{include/JtagScanDiagram.tex}
	\caption[{One four-bit data-register scan.}]{One four-bit data-register scan. The state row is the sequence the port walks under the \pin{TMS} pattern above it. Data is shifted least-significant bit first, one bit per \pin{TCK}; \pin{TDI} is sampled on the rising edge and \pin{TDO} changes on the falling one, which is why the two rows are offset by half a cycle. A real scan is 32 or 41 bits wide rather than four, but the shape is exactly this.}
	\label{fig:jtag-scan}
\end{figure}

The instruction register selects \textit{which} data register a subsequent DR
scan will reach. Every JTAG device is therefore fully described by three things:
the width of its instruction register, the set of instruction values it
recognises, and the width and meaning of the data register each one selects.\ifdebugenable{}
For this chip, that is Table \ref{t:debug-ir}.\fi

Two consequences are worth stating because they surprise people. First, the port
is entirely passive with respect to the chip's own clocks: it will answer while
the processor is running, while it is stopped, and while it is held in reset.
Second, because the standard defines a \textbf{BYPASS} instruction that selects a
single flip-flop, several devices can be wired in a chain --- \pin{TDO} of one to
\pin{TDI} of the next --- and a debugger can talk to any one of them while the
others contribute one bit each of padding.

\subsection{The RISC-V Debug Stack} \label{ss:debug-stack}

JTAG by itself only gives you a way to shift bits in and out. The RISC-V debug
specification defines what those bits \textit{mean}, and it does so in four
layers. Each layer knows nothing about the internals of the one below it, which
is what allows the same debugger software to drive very different chips.

\begin{center}
\begin{tabular}{ l p{11.0cm} }
	\textbf{Layer} & \textbf{What it is} \\ \hline
	\textbf{DTM} & The \textit{Debug Transport Module}. It is the JTAG-facing half: the TAP itself, plus the logic that turns one DR scan into one transaction on the interface below. Its job is transport and nothing else --- it has no idea what a hart is. \\
	\rowcolor{tablehighlightcolor} \textbf{DMI} & The \textit{Debug Module Interface}. A tiny address/data bus --- here 7 address bits and 32 data bits --- with three operations: no-op, read, write. This is the seam. Anything that can drive a DMI can debug the chip; JTAG is merely the transport this chip provides. \\
	\textbf{DM} & The \textit{Debug Module}. A block of registers sitting on the DMI that implements \textit{run control}: select a hart, halt it, ask what state it is in, run a small program on it, resume it. There is one DM for the whole chip and it can reach every hart. \\
	\rowcolor{tablehighlightcolor} \textbf{Debug mode} & A state of the \textit{hart itself} --- effectively a privilege level above machine mode. A hart in debug mode has stopped executing your program, is executing the debugger's code instead, and has a small set of control registers that only exist while it is there. \\ \hline
\end{tabular}
\end{center}

The layer that most often confuses newcomers is the last one. \textbf{A halted
hart is not a frozen hart.} On this architecture, halting a processor does not
stop its clock and freeze it mid-instruction; it \textit{diverts} it, between
instructions, to a fixed address where the debugger's own code is waiting. The
hart then sits in a loop at that address, still fetching and executing, waiting
to be told what to do. When the debugger wants to read a register, it does not
reach into the register file with wires --- it writes a two- or three-instruction
program into memory, tells the hart to run it, and reads back where the program
put the answer.

This is the single most important thing to understand about RISC-V debug, and
almost every property in the rest of this chapter follows from it: that the
debugger needs somewhere in memory to put code, that reading a register can
report an \textit{exception}, that the hart's own \asminline{s0} and
\asminline{s1} registers need saving before anything else can happen, and that
memory accesses on behalf of the debugger go through the ordinary bus, with
ordinary arbitration, exactly as the program's own would.

Figure \ref{fig:debug-stack} shows those layers as they are actually built, in a
configuration that includes the debug system: which clock each part runs on,
what crosses between them, and the three ways the Debug Module reaches the rest
of the chip.

% Generated by LatexUserGuide.GenerateDebugStackDiagram (make chip).
\begin{figure}[htpb]
	\centering
	\resizebox{\linewidth}{!}{\input{include/DebugStackDiagram.tex}}
	\caption[{The debug path, from the probe to a hart.}]{The debug path, from the probe to a hart, in a configuration built with debug support. The transport and the Debug Module sit in \textit{different clock domains}, and the wall between them is crossed by exactly two signals --- one toggle each way, each carrying a payload that is held still while its toggle is in flight. Everything to the right of that wall is ordinary chip fabric, and the one Debug Module reaches it in three ways: it halts and resumes the harts over direct wires, it reads and writes memory as one more master on the same arbiter the harts use, and it plants its own trampoline into the debug program page. The count of hart tiles follows this configuration.}
	\label{fig:debug-stack}
\end{figure}

\ifdebugenable

\subsection{The Debug Port} \label{ss:debug-port}

The debug port occupies five pins, which are present only on the LQFP-100
package; a configuration built for a smaller package has the debug system in
silicon but nowhere to bond it.

\begin{center}
\begin{tabular}{ l l l p{6.6cm} }
	\textbf{Pin} & \textbf{Signal} & \textbf{I/O} & \textbf{Idle state when unconnected} \\ \hline
	47 & \pin{TCK} & Input & Pulled \textbf{low}. \pin{TCK} idles low. \\
	\rowcolor{tablehighlightcolor} 48 & \pin{TMS} & Input & Pulled \textbf{high}. This is what IEEE 1149.1 requires, and it is what allows five clocks from a disconnected probe to still return the port to Test-Logic-Reset. \\
	49 & \pin{TDI} & Input & Pulled \textbf{high}. \\
	\rowcolor{tablehighlightcolor} 50 & \pin{TDO} & Output & Driven. Held quiet between shift states rather than toggling, so the pin is well behaved in a multi-device chain. \\
	51 & \pin{TRSTn} & Input & Pulled \textbf{low} --- see below. \\ \hline
\end{tabular}
\end{center}

\textbf{The pull on \pin{TRSTn} is a deliberate, documented deviation from IEEE
1149.1.} The standard wants \pin{TRSTn} to idle high, so that a port with a
floating reset pin is available to a probe that turns up later. This chip pulls
it \textit{low}, which holds the TAP in reset. The reasoning is that the debug
transport here is optional: a board that does not use it leaves all five pins
unconnected, and in that case the desirable behaviour is a port that is
completely inert --- issuing nothing, clocking nothing, and unable to reach the
Debug Module at all --- rather than a port that is quietly attachable by anyone
who touches the pins. \textbf{Fail-safe-inert was chosen over
fail-safe-attachable.} A board that wants to use the port must therefore drive
\pin{TRSTn} high, and cannot rely on the pull to do it.

The five pins are ordinary digital I/O on the 3.3~V \texttt{VDDPST}/\texttt{VSSPST}
domain and may be left unconnected; the chip is unaffected. They are deliberately
placed on the south and east edges, away from the analog band on the north edge.

\textbf{Identification.} A DR scan of the IDCODE register returns a 32-bit value
that identifies the part before the debugger knows anything else about it:

\begin{center}
\begin{tabular}{ l l p{8.0cm} }
	\textbf{Chip} & \textbf{IDCODE} & \textbf{Decomposition} \\ \hline
	Castalia & \texttt{0x1CA57EEF} & version \texttt{1}, part number \texttt{0xCA57}, manufacturer \texttt{0x777}, mandatory LSB \texttt{1} \\
	\rowcolor{tablehighlightcolor} Argus & \texttt{0x1A265EEF} & version \texttt{1}, part number \texttt{0xA265}, manufacturer \texttt{0x777}, mandatory LSB \texttt{1} \\ \hline
\end{tabular}
\end{center}

The manufacturer field is \texttt{0x777}, which is \textit{not} a valid JEDEC
JEP106 code and is not registered to anyone. That is intentional and it is the
honest answer for a non-commercial teaching chip: rather than borrow an
identifier belonging to a real company, the field carries a value that is
recognisably invalid. Discovery software should identify this part by its full
IDCODE, not by its manufacturer field.

\subsection{The Debug Transport Module} \label{ss:debug-dtm}

The DTM implements the TAP described in Section \ref{ss:debug-jtag} and four
data registers. Its instruction register is \textbf{five bits} wide and resets
to the IDCODE instruction.

\begin{table}[htb]
	\centering
	\begin{tabular}{ l l l p{6.4cm} }
	\caption{Debug Transport Module instruction register values} \label{t:debug-ir}
	\textbf{IR value} & \textbf{Instruction} & \textbf{DR width} & \textbf{Selects} \\ \hline
	\texttt{0x01} & IDCODE & 32 & The identification value above. This is the reset instruction, so the first DR scan after a port reset always returns the IDCODE. \\
	\rowcolor{tablehighlightcolor} \texttt{0x10} & DTMCS & 32 & The transport's own control and status register. \\
	\texttt{0x11} & DMI & 41 & One transaction on the Debug Module Interface. \\
	\rowcolor{tablehighlightcolor} \texttt{0x1F} & BYPASS & 1 & A single flip-flop, for use in a multi-device chain. \\
	\textit{anything else} & --- & 1 & \textbf{BYPASS.} Every unrecognised instruction selects the one-bit chain, so the length of the selected register never depends on what was scanned before. \\
	\hline
	\end{tabular}
\end{table}

\textbf{Test-Logic-Reset sets the instruction register to IDCODE and does
nothing else.} In particular it does \textit{not} clear the sticky status
described below, and does not discard a result waiting to be read. That is what
the \bitfield{dmireset} strobe is for, and it is why a debugger that has just
recovered a lost TAP still has to clear the status explicitly. \pin{TRSTn}, by
contrast, resets the entire transport.

\textbf{DTMCS} reports what the transport can do and carries the two recovery
strobes:

\begin{center}
\begin{tabular}{ l l p{8.6cm} }
	\textbf{Bits} & \textbf{Field} & \textbf{Meaning} \\ \hline
	\texttt{[3:0]} & \bitfield{version} & \texttt{1} --- this transport satisfies both the 0.13 and the 1.0 debug specifications. \\
	\rowcolor{tablehighlightcolor} \texttt{[9:4]} & \bitfield{abits} & \texttt{7} --- the DMI address is seven bits wide, so the DMI data register is $7 + 32 + 2 = 41$ bits. \\
	\texttt{[11:10]} & \bitfield{dmistat} & Sticky transaction status; see below. \\
	\rowcolor{tablehighlightcolor} \texttt{[14:12]} & \bitfield{idle} & \texttt{7} --- the number of Run-Test/Idle clocks the debugger should insert between a DMI scan and the scan that collects its result. \\
	\texttt{[16]} & \bitfield{dmireset} & Write \texttt{1} to clear \bitfield{dmistat} and abandon any outstanding transaction. Reads as \texttt{0}. \\
	\rowcolor{tablehighlightcolor} \texttt{[17]} & \bitfield{dmihardreset} & Write \texttt{1} to do the above \textit{and} clear the stored result. Reads as \texttt{0}. \\ \hline
\end{tabular}
\end{center}

\textbf{The \bitfield{idle} value is truthful, and that is a deviation worth
knowing about.} Many implementations advertise \texttt{0} here while in fact
requiring idle clocks, leaving the debugger to discover the requirement by
failing. The value \texttt{7} reported here is derived from the real latency of
the slowest transaction class. A DMI access to one of the Debug Module's own
registers completes in eight to ten cycles of the internal main clock; an access
to \register{data0} or a program-buffer word is \textit{proxied} into shared
memory through the arbiter and takes tens of cycles. Against an assumed
\pin{TCK} period of at least 140~ns --- that is, \pin{TCK} at or below
7.14~MHz --- against a 24~MHz main clock, the worst class needs six idle cycles,
and \texttt{7} gives one cycle of margin. It is also the largest value a
three-bit field can express.

\textbf{\bitfield{idle} is a recommendation, not a guarantee.} At a faster
\pin{TCK}, or while the shared bus is heavily contended by the harts, a
transaction can still be outstanding when the debugger comes back for it. That
is not an error; it is reported as \textit{busy} and the correct response is to
clear the status and reissue.

\textbf{The DMI data register} is 41 bits, laid out as in Figure
\ref{fig:dmi-field}: \bitfield{op}\texttt{[1:0]}, \bitfield{data}\texttt{[33:2]},
\bitfield{address}\texttt{[40:34]}.

% Generated by LatexUserGuide.GenerateDmiFieldDiagram (make chip).
\begin{figure}[htpb]
	\centering
	\input{include/DmiFieldDiagram.tex}
	\caption[{The 41-bit DMI data register.}]{The 41-bit DMI data register --- one scan of it is one transaction on the Debug Module Interface.}
	\label{fig:dmi-field}
\end{figure} On the way \textit{in}, \bitfield{op} is
\texttt{00} for no-operation, \texttt{01} for read and \texttt{10} for write. On
the way \textit{out}, it is \texttt{00} for success, \texttt{10} for failed and
\texttt{11} for busy. \textbf{A no-operation scan never starts anything}, which
is what makes a capture-only scan --- shifting in zeros purely to read out the
previous result --- safe.

\textbf{\bitfield{dmistat} is sticky, and it has two flavours.} It takes the
value \texttt{3} (busy) if a scan is captured while a transaction is still in
flight, or if a new transaction is scanned in while one is outstanding; and
\texttt{2} (failed) if a transaction completes with a failure --- which happens,
for instance, on a DMI access to an address the Debug Module does not implement.
\textbf{While \bitfield{dmistat} is nonzero, in either flavour, further DMI
scans are dropped.} This is deliberate: it stops a debugger from issuing a long
sequence of writes into the void after the first one failed. The recovery is
always the same --- write \bitfield{dmireset}, then reissue --- and the status is
cleared by \bitfield{dmireset} or \bitfield{dmihardreset} and by nothing else.
Neither Test-Logic-Reset nor a system reset clears it.

Figure \ref{fig:dmi-crossing} follows a single transaction across the clock
boundary, and is where the \bitfield{idle} value above comes from.

% Generated by LatexUserGuide.GenerateDmiCrossingDiagram (make chip).
\begin{figure}[htpb]
	\centering
	\resizebox{\linewidth}{!}{\input{include/DmiCrossingDiagram.tex}}
	\caption[{One DMI transaction crossing between the two clock domains.}]{One DMI transaction crossing between the two clock domains, drawn on the internal clock. The two events in the debug probe's own clock domain are marked above the waveforms rather than drawn, because the ratio between the two clocks is a board-level choice. The shaded cycles are the request handshake: the request is presented for exactly two cycles and retired as soon as the acceptance is observed, and the note below explains why holding it any longer would make the same request be answered twice.}
	\label{fig:dmi-crossing}
\end{figure}

\subsection{The Debug Module} \label{ss:debug-dm}

There is one Debug Module for the whole chip. It sits on the always-running main
clock, it is reachable by the DTM through the DMI, and --- because it must be
able to place instructions where a hart will fetch them --- it is also a
\textit{master} on the shared-memory arbiter, alongside the harts (Section
\ref{s:multicore}). Table \ref{t:debug-dmi} is its DMI address map.

\begin{center}
\begin{longtable}[c]{ l l p{9.0cm} }
\caption{Debug Module registers, at their DMI addresses} \label{t:debug-dmi} \\
\hline \textbf{DMI} & \textbf{Register} & \textbf{Purpose} \\ \hline \endfirsthead
\multicolumn{3}{c}{\textit{\tablename\ \thetable\ continued from previous page}} \\ \hline
\textbf{DMI} & \textbf{Register} & \textbf{Purpose} \\ \hline \endhead
\hline \multicolumn{3}{c}{\textit{\tablename\ \thetable\ continued on next page}} \\ \endfoot \hline \endlastfoot
\texttt{0x04} & \register{data0} & The single abstract-command data word. Backed by shared memory, not by a flip-flop. \\
\rowcolor{tablehighlightcolor} \texttt{0x10} & \register{dmcontrol} & Run control: enable the module, select a hart, halt it, resume it. \\
\texttt{0x11} & \register{dmstatus} & The selected hart's state, and the module's capabilities. Read-only. \\
\rowcolor{tablehighlightcolor} \texttt{0x12} & \register{hartinfo} & Where \register{data0} lives, so a debugger can find it. Read-only. \\
\texttt{0x13} & \register{haltsum1} & Reads zero. Present, and answering successfully, because debuggers probe it. \\
\rowcolor{tablehighlightcolor} \texttt{0x16} & \register{abstractcs} & Abstract-command status: busy, error code, and the module's sizes. \\
\texttt{0x17} & \register{command} & Write to start an abstract command. Reads zero. \\
\rowcolor{tablehighlightcolor} \texttt{0x18} & \register{abstractauto} & Reads zero; automatic command re-execution is not implemented. \\
\texttt{0x20} & \register{progbuf0} & Program-buffer word 0. Backed by shared memory. \\
\rowcolor{tablehighlightcolor} \texttt{0x21} & \register{progbuf1} & Program-buffer word 1. Backed by shared memory. \\
\texttt{0x32} & \register{dmcs2} & Halt-group membership of the selected hart. \\
\rowcolor{tablehighlightcolor} \texttt{0x40} & \register{haltsum0} & One bit per hart: halted and reachable. The only status register that ignores the hart selection. \\
\hline
\end{longtable}
\end{center}

\textbf{Any other DMI address returns \textit{failed}}, which --- because
\bitfield{dmistat} is sticky --- will also stall the transport until the debugger
clears it. Probing for unimplemented registers is therefore not free, and this is
why \register{haltsum1} and \register{abstractauto} answer with a successful zero
instead of a failure: they are addresses debuggers routinely read during
discovery, and failing them would punish a debugger for being thorough.

\textbf{\register{dmcontrol}} is the register everything starts with.
\bitfield{dmactive} (bit 0) must be written \texttt{1} before the module will do
anything; writing it \texttt{0} resets the module's own state --- but
\textit{not} the harts. \bitfield{hartsel} (bits \texttt{[25:16]}) selects which
hart the status and command registers refer to; all ten bits are implemented, and
selecting a hart index that does not exist is reported honestly as
\textit{nonexistent} rather than silently clamped to the highest real hart, which
is what makes hart-count discovery work. \bitfield{haltreq} (bit 31) is a level:
while it is set, the selected hart is asked to halt. \bitfield{resumereq} (bit 30)
is a write-one request to resume; it is queued rather than refused, and it is
ignored for a hart that is not halted or is not reachable.
\bitfield{setresethaltreq} and \bitfield{clrresethaltreq} (bits 3 and 2) arm and
disarm halt-on-reset for the selected hart. \bitfield{ndmreset},
\bitfield{hartreset} and \bitfield{hasel} are not implemented; they read zero and
ignore writes.

\textbf{\register{dmstatus}} reports \bitfield{version}~$=3$ (debug specification
1.0), \bitfield{authenticated}~$=1$ (there is no authentication on this chip),
\bitfield{hasresethaltreq}~$=1$ and \bitfield{impebreak}~$=1$. Its state is
reported through the usual paired bits --- \bitfield{allhalted}/\bitfield{anyhalted},
\bitfield{allrunning}/\bitfield{anyrunning}, \bitfield{allunavail}/\bitfield{anyunavail},
\bitfield{allnonexistent}/\bitfield{anynonexistent},
\bitfield{allresumeack}/\bitfield{anyresumeack},
\bitfield{allhavereset}/\bitfield{anyhavereset}. Because only one hart is ever
selected, both bits of every pair always carry the same value.

\textbf{\register{abstractcs}} reports \bitfield{progbufsize}~$=2$ and
\bitfield{datacount}~$=1$: two program-buffer words and one data word.
\bitfield{busy} (bit 12) is set while a command is running --- and the module
remains readable while it is, so the debugger can poll. \bitfield{cmderr} (bits
\texttt{[10:8]}) is the error code, and it is \textbf{write-one-to-clear}:

\begin{center}
\begin{tabular}{ l l p{9.2cm} }
	\textbf{\bitfield{cmderr}} & \textbf{Meaning} & \textbf{Raised when} \\ \hline
	\texttt{0} & none & --- \\
	\rowcolor{tablehighlightcolor} \texttt{1} & busy & A command was written while one was running, or \register{data0}/a program-buffer word was accessed while the module's bus engine was busy. \\
	\texttt{2} & not supported & The command type is not access-register, or the transfer size is not 32-bit. \\
	\rowcolor{tablehighlightcolor} \texttt{3} & exception & The hart took an exception while executing the command. \\
	\texttt{4} & halt/resume & The selected hart was not halted when the command was written. \\
	\rowcolor{tablehighlightcolor} \texttt{7} & other & A command did not complete within the module's internal timeout. \\ \hline
\end{tabular}
\end{center}

\textbf{A command written while \bitfield{cmderr} is already nonzero is ignored
entirely} --- it does not run, and it does not clear the error on its way past.
Software must clear \bitfield{cmderr} before issuing the next command.

\textbf{\register{command}} implements one command type: \textbf{access
register}. It transfers one 32-bit value between \register{data0} and a register
of the halted hart, optionally executing the program buffer afterwards. The
fields are \bitfield{aarsize} (\texttt{[22:20]}, which must be \texttt{2} for
32-bit), \bitfield{postexec} (bit 18), \bitfield{transfer} (bit 17),
\bitfield{write} (bit 16) and \bitfield{regno} (\texttt{[15:0]}). Register
numbers \texttt{0x1000}--\texttt{0x101F} are the general-purpose registers
\asminline{x0}--\asminline{x31}; any other number is treated as a control and
status register address. \textbf{Access memory is deliberately not
implemented} and returns \textit{not supported}. Memory is reached instead by
writing an \asminline{lw} or \asminline{sw} into the program buffer and executing
it on the halted hart --- see Section \ref{ss:debug-limits}.

\textbf{\register{data0} and the program-buffer words are memory, not
registers.} A DMI access to \texttt{0x04}, \texttt{0x20} or \texttt{0x21}
becomes a read or write of a word in shared memory, issued by the Debug Module
through the ordinary arbiter. This is why those three addresses are the slow
class that sets \bitfield{idle} to \texttt{7}, and why an access to them while
the module's bus engine is busy answers successfully with zero data \textit{and}
sets \bitfield{cmderr} to \textit{busy} --- the answer arrives, but it is not the
answer you wanted, and the error code is how you find out.

The program buffer is two DMI-visible words plus an \textbf{implicit third
word}, which is why \bitfield{impebreak} reads \texttt{1}. Software may therefore
place up to three instructions' worth of work in the buffer, with the third slot
supplied by hardware.

\subsection{Debug Mode in the Hart} \label{ss:debug-hart}

A hart enters debug mode for one of four reasons, recorded in
\register{dcsr}.\bitfield{cause}:

\begin{center}
\begin{tabular}{ l l p{9.0cm} }
	\textbf{\bitfield{cause}} & \textbf{Reason} & \textbf{Notes} \\ \hline
	\texttt{1} & \asminline{EBREAK} & A software breakpoint, taken only when \register{dcsr}.\bitfield{ebreakm} is set (or when the hart is already in debug mode). \\
	\rowcolor{tablehighlightcolor} \texttt{3} & Halt request & The Debug Module asserted \bitfield{haltreq}. \\
	\texttt{4} & Single step & The one instruction permitted by \register{dcsr}.\bitfield{step} has retired. \\
	\rowcolor{tablehighlightcolor} \texttt{5} & Halt on reset & \bitfield{setresethaltreq} was armed when this hart came out of reset. \\ \hline
\end{tabular}
\end{center}

On entry, the hart saves its program counter to \register{dpc}, records the
cause and its previous privilege level in \register{dcsr}, sets debug mode, and
jumps to a fixed address --- \textbf{\texttt{0x00010780}}, in the shared memory
window. It is now executing the debugger's code.

Four control registers exist while it is there, and \textbf{only} while it is
there:

\begin{center}
\begin{tabular}{ l l p{9.2cm} }
	\textbf{Address} & \textbf{Register} & \textbf{Contents} \\ \hline
	\texttt{0x7B0} & \register{dcsr} & Debug control and status: \bitfield{xdebugver} \texttt{[31:28]} reads \texttt{4}; \bitfield{ebreakm} (15) makes \asminline{EBREAK} enter debug mode; \bitfield{ebreaku} (12) does the same for user mode, where user mode is configured; \bitfield{cause} \texttt{[8:6]} is read-only and set by hardware; \bitfield{step} (2) enables single-stepping; \bitfield{prv} \texttt{[1:0]} is the privilege level to return to. \\
	\rowcolor{tablehighlightcolor} \texttt{0x7B1} & \register{dpc} & The saved program counter. Bit 0 is hardwired zero; bit 1 is writable only where the compressed extension is configured. \\
	\texttt{0x7B2} & \register{dscratch0} & A 32-bit scratch word for the debugger's own use. \\
	\rowcolor{tablehighlightcolor} \texttt{0x7B3} & \register{dscratch1} & A second scratch word. \\ \hline
\end{tabular}
\end{center}

\textbf{Accessing any of them from machine or user mode is an illegal
instruction, and the denial is wider than the four addresses.} Every CSR
instruction naming an address in \texttt{0x7B0}--\texttt{0x7BF} is rejected
outside debug mode. This is not an optimisation: it means a program cannot
accidentally --- or deliberately --- read the debugger's saved state or forge an
entry.

\textbf{\asminline{DRET} (\texttt{0x7B200073}) leaves debug mode}, restoring the
program counter from \register{dpc} and the privilege level from
\register{dcsr}.\bitfield{prv}. Outside debug mode it is an illegal instruction.

Five behaviours are worth calling out because they are the ones that decide
whether a debug session works:

\begin{itemize}
	\item \textbf{The halt request is unmaskable.} It is recognised at fourteen points in the processor's instruction sequencer --- the thirteen at which an interrupt is recognised, \textit{plus the terminal trap state}. A hart that has taken an illegal instruction and stopped retiring, which is otherwise unrecoverable short of reset (Section \ref{ss:priv-legacy}), can therefore be halted and inspected. Neither \register{mstatus}.\bitfield{MIE} nor \register{mtrapctl}.\bitfield{LEGACY} has any effect on it. Note that this rescues a wedged hart only if something is there to request the halt: with no debugger attached the wedge is exactly as terminal as it always was.
	\item \textbf{An \asminline{EBREAK} records its own address in \register{dpc}}, not the address after it. A debugger resuming past a software breakpoint must add the instruction's length itself. Every other cause is taken between instructions, so \register{dpc} is already the resume address.
	\item \textbf{No interrupt is taken in debug mode.} Interrupt sources remain pending as levels and are delivered after the \asminline{DRET}, so a debug session neither loses interrupts nor lets the application's handlers run underneath it.
	\item \textbf{An exception taken in debug mode re-enters debug mode} at the same fixed address, leaving \register{dpc} and \register{dcsr} untouched, and without disturbing \register{mepc}, \register{mcause} or \register{mtval}. This is what makes a faulting abstract command --- reading a CSR the configuration does not implement, say --- report \bitfield{cmderr}~$=3$ instead of destroying the session.
	\item \textbf{Halt-on-reset is sampled once, at reset release.} Arming \bitfield{setresethaltreq} affects the \textit{next} time that hart comes out of reset; it will not halt a hart that is already running. This is what lets a debugger catch a hart before its first instruction, and it is why arming it does not disturb a running system.
\end{itemize}

Finally: \textbf{single-step diverts at the stepped instruction's own dispatch},
so exactly one instruction retires per step. Stepping is inactive while in debug
mode, so the debugger's own stub does not step itself.

Figure \ref{fig:debug-mode-state} collects the whole of the above into the two
states a hart can be in.

% Generated by LatexUserGuide.GenerateDebugModeStateDiagram (make chip).
\begin{figure}[htpb]
	\centering
	\resizebox{\linewidth}{!}{\input{include/DebugModeStateDiagram.tex}}
	\caption[{Debug mode, from the hart's point of view.}]{Debug mode, from the hart's point of view. Debug mode behaves as a privilege level above machine mode: the hart has not stopped, it has been diverted to code the debugger controls, and it returns by executing \asminline{dret}. The self-loop is the property that makes a faulting abstract command recoverable rather than fatal.}
	\label{fig:debug-mode-state}
\end{figure}

\subsection{Halt Groups and Unavailable Harts} \label{ss:debug-groups}

Two facilities exist because this is a multi-hart chip.

\textbf{Halt groups} solve the problem that stopping one hart of a
cooperating set, while the others run on, changes the very behaviour you were
trying to observe. \register{dmcs2} assigns the selected hart to one of eight
groups; group \texttt{0}, the reset value, means ``no group''. \textbf{When any
member of a group halts --- for any reason, including its own
\asminline{EBREAK} --- the Debug Module requests a halt of every other member.}
Setting a breakpoint in one hart therefore stops the whole group at very nearly
the same instant. Note that this is halt-only: resuming a group is done by
resuming its members individually.

\textbf{Unavailable is a third state, and it is not the same as halted.} This
chip power-gates the harts (Section \ref{s:multicore}); a hart whose power domain
is off cannot answer the Debug Module at all, and the signals that would report
its state are clamped to zero at the domain boundary. A clamped ``not halted'' is
indistinguishable from ``running'', so the Debug Module does not rely on it: it
takes a separate, deliberate \textit{unavailable} signal from the power
controller, and consults it \textbf{before} halted or running. The reported state
is therefore resolved in strict order --- \textbf{nonexistent, then unavailable,
then halted, then running} --- and a gated hart reports
\bitfield{allunavail}/\bitfield{anyunavail} with both halted and running clear.
It is never misreported as running.

Two consequences follow. First, a halt request aimed at a gated hart does
nothing; the debugger must power the hart up first, through the power controller,
exactly as software would. Second, \textbf{a hart returning from power-down
cold-boots the boot ROM and parks --- it does not resume where it left off}. If
you want to catch it as it comes back, arm halt-on-reset before powering it up.
Note that hart 0 is not exempt from any of this.

\subsection{The Debug Program Page} \label{ss:debug-page}

Because a halted hart executes code, the debug system needs memory it owns. It
uses one 512-byte region of the shared window, listed here so that application
software does not collide with it.

\begin{center}
\begin{tabular}{ l p{9.6cm} }
	\textbf{Address} & \textbf{Contents} \\ \hline
	\texttt{0x10680} & \register{data0} --- the abstract-command data word. This is the address \register{hartinfo} reports. \\
	\rowcolor{tablehighlightcolor} \texttt{0x10684} & \register{progbuf0} \\
	\texttt{0x10688} & \register{progbuf1} \\
	\rowcolor{tablehighlightcolor} \texttt{0x1068C} & The implicit third program-buffer word, written by the Debug Module. \\
	\texttt{0x10690}--\texttt{0x106EC} & Reserved for the Debug Module. \\
	\rowcolor{tablehighlightcolor} \texttt{0x106F0}, \texttt{0x106F4} & The halted session's saved \asminline{s0} and \asminline{s1}. \\
	\texttt{0x10700}\,$+$\,$4h$ & The per-hart handshake word between the Debug Module and hart $h$. \\
	\rowcolor{tablehighlightcolor} \texttt{0x10780}--\texttt{0x1087F} & The debug entry page, 64 words: the entry stub, the area into which the Debug Module writes each abstract command, and the common exit sequence. \\ \hline
\end{tabular}
\end{center}

Figure \ref{fig:debug-page} draws the same region to scale, with the one
boundary inside it that software has to know about.

% Generated by LatexUserGuide.GenerateDebugPageDiagram (make chip).
\begin{figure}[htpb]
	\centering
	\resizebox{\linewidth}{!}{\input{include/DebugPageDiagram.tex}}
	\caption[{The debug program page.}]{The debug program page. The boot ROM zero-fills shared memory only up to \texttt{0x107FF}, so the dashed line marks where that guarantee stops: everything above it is code the Debug Module writes before anything can read it, which is why the page needs no staging and cannot be read-only memory.}
	\label{fig:debug-page}
\end{figure}

\textbf{The whole region \texttt{0x10680}--\texttt{0x1087F} is reserved.}
Application software, including software on harts that are not being debugged,
must not write to it in a configuration built with debug support. A stray store
into the entry page does not merely corrupt data --- it corrupts the code a
halted hart is executing.

Two details matter to anyone writing code that runs \textit{in} that page. It is
assembled without compressed instructions, because execution of compressed or
misaligned instructions from the shared window is outside the tested envelope
(Section \ref{s:multicore}). And the debug system's use of \asminline{s0} and
\asminline{s1} is why those two registers, uniquely, are served from their saved
copies when a debugger reads them: you see the interrupted program's values, not
the stub's working values.

\subsection{A Worked Halt, Read and Resume} \label{ss:debug-worked}

The following is a complete session at the DMI level: attach, halt hart 2, read
its \asminline{a1} register, and resume it. Each numbered step is one or more
JTAG scans. A ``DMI read'' is two scans --- one to issue the read, and one, after
the recommended idle clocks, to collect the result --- because the transport
returns the \textit{previous} transaction's result on every capture.

\begin{enumerate}
	\item \textbf{Reset the port and identify the chip.} Drive \pin{TRSTn} high, then clock \pin{TMS} high for five \pin{TCK} cycles to reach Test-Logic-Reset. The instruction register is now IDCODE, so a 32-bit DR scan returns \texttt{0x1CA57EEF}.
	\item \textbf{Read DTMCS.} Scan \texttt{0x10} into the instruction register, then a 32-bit DR scan. Confirm \bitfield{version}~$=1$ and \bitfield{abits}~$=7$, and note \bitfield{idle}~$=7$: insert seven Run-Test/Idle clocks after every DMI scan from here on.
	\item \textbf{Select the DMI register.} Scan \texttt{0x11} into the instruction register. Every scan below is a 41-bit DR scan carrying $\{$\bitfield{address}, \bitfield{data}, \bitfield{op}$\}$.
	\item \textbf{Enable the Debug Module.} Write \texttt{0x00000001} to \register{dmcontrol} (\texttt{0x10}): $\{$\texttt{0x10}, \texttt{0x00000001}, write$\}$.
	\item \textbf{Select hart 2 and request a halt.} Write \texttt{0x80020001} to \register{dmcontrol} --- \bitfield{haltreq} set (bit 31), \bitfield{hartsel}~$=2$ (bits \texttt{[25:16]}), \bitfield{dmactive} set.
	\item \textbf{Wait for the halt.} Read \register{dmstatus} (\texttt{0x11}) repeatedly until \bitfield{allhalted} (bit 9) is set. If \bitfield{allunavail} is set instead, hart 2 is powered down; power it up before continuing. If \bitfield{allnonexistent} is set, this configuration has fewer than three harts.
	\item \textbf{Drop the halt request.} Write \texttt{0x00020001} to \register{dmcontrol}, leaving hart 2 selected. The hart stays halted --- halting is a state of the hart, not a level the Debug Module holds.
	\item \textbf{Clear any stale error.} Write \texttt{0x00000700} to \register{abstractcs} (\texttt{0x16}) to clear \bitfield{cmderr}. A command written while an error stands is ignored.
	\item \textbf{Read \asminline{a1}.} \asminline{a1} is \asminline{x11}, so \bitfield{regno}~$=$~\texttt{0x100B}. Write \texttt{0x0022100B} to \register{command} (\texttt{0x17}): \bitfield{aarsize}~$=2$ (\texttt{0x00200000}) $+$ \bitfield{transfer} (\texttt{0x00020000}) $+$ \bitfield{regno}.
	\item \textbf{Wait for the command.} Read \register{abstractcs} until \bitfield{busy} (bit 12) is clear, then check \bitfield{cmderr} \texttt{[10:8]}. Zero means the value is ready.
	\item \textbf{Collect the value.} Read \register{data0} (\texttt{0x04}). This one is proxied through shared memory, so it is the transaction the seven idle clocks were sized for.
	\item \textbf{Resume.} Write \texttt{0x40020001} to \register{dmcontrol} --- \bitfield{resumereq} (bit 30), hart 2 still selected. Poll \register{dmstatus} until \bitfield{allresumeack} (bit 17) is set.
\end{enumerate}

% Generated by LatexUserGuide.GenerateDebugSwimlaneDiagram (make chip).
\begin{figure}[htpb]
	\centering
	\resizebox{\linewidth}{!}{\input{include/DebugSwimlaneDiagram.tex}}
	\caption[{The twelve steps, across the four agents that perform them.}]{The twelve steps above, across the four agents that perform them. The numbers are the list items. Note how little of the sequence is the debugger talking to the hart directly: most of it is the debugger arranging for the \textit{hart} to do the work, which is the consequence of there being no independent path from the debug port to memory.}
	\label{fig:debug-swimlane}
\end{figure}

To write a register instead of reading one, put the value in \register{data0}
first and set \bitfield{write} (bit 16) in the command. \textbf{To read or write
\textit{memory}}, put an \asminline{lw} or \asminline{sw} in the program buffer,
put the address in a register of the halted hart, and issue a command with
\bitfield{postexec} set; the hart executes your instruction with its own
privilege and its own view of the address space.

If a scan ever returns \bitfield{op}~$=$~\texttt{10} or \texttt{11}, the
transport has latched a sticky status and is dropping further scans. Read
DTMCS, write \bitfield{dmireset}, and reissue.

\subsection{Limitations and Status} \label{ss:debug-limits}

This section states plainly what the debug system does not do, and what is not
yet finished.

\begin{itemize}
	\item \textbf{There is no System Bus Access, and there never will be.} Many RISC-V Debug Modules include a second bus master that reads and writes memory directly, independently of any hart. This one does not, by design. Memory is reached only by executing \asminline{lw} and \asminline{sw} on a halted hart through the program buffer. The consequences are worth understanding: memory accesses on behalf of the debugger go through the same arbiter, with the same permissions and the same view of the address space, as the hart's own --- which is a feature, not a workaround --- but they require a hart that is halted and healthy, so a chip with no hart in a debuggable state cannot have its memory inspected.
	\item \textbf{There are no hardware triggers.} Address and data watchpoints, and hardware breakpoints on instruction fetch, are not implemented; the trigger registers are absent. Breakpoints are software breakpoints: replace the instruction with \asminline{EBREAK} and set \register{dcsr}.\bitfield{ebreakm}. This means breakpoints cannot be set in the boot ROM or in any read-only region.
	\item \textbf{The debug system requires the standard trap architecture.} \texttt{debug.enable} cannot be configured without \texttt{priv.trapCsr} (Section \ref{s:privarch}); the generator rejects the combination. \asminline{EBREAK} is decoded on the standard privileged instruction path, so a chip without it could not recognise a software breakpoint at all.
	\item \textbf{The debug system is a configuration option and is off by default.} A build without \texttt{debug.enable} contains none of this hardware --- no debug port, no Debug Module, no debug mode --- and its processor is bit-identical to one built before the debug system existed. See Table \ref{t:chip-config} for the setting in this build.
	\item \textbf{\bitfield{ndmreset} and \bitfield{hartreset} are not implemented.} A debugger cannot reset the system or an individual hart through the Debug Module; they read zero and ignore writes. Resetting the chip is done with the \pin{resetn} pin (Section \ref{s:startup}).
	\item \textbf{Resume groups are not implemented.} Halt groups stop a set of harts together; resuming them is done one at a time.
	\item \textbf{The entry page is volatile, and a wrong first word wedges the hart that halts into it.} The code a hart runs when it enters debug mode occupies \texttt{0x10780} onwards, and the Debug Module places it there itself --- once when a debugger raises \bitfield{dmactive}, and again before it acts on any hart it has just seen halt. A production part therefore needs nothing staged in advance, and there is no debug ROM: that page \textit{cannot} be read-only memory, because the Debug Module rewrites part of it at every abstract command. What it is instead is ordinary shared RAM, so any master on the chip can overwrite it. Most such damage repairs itself, because the next halt rewrites the page before the Debug Module depends on it --- but not damage to the \textbf{first} word. A hart that halts into a wrong first word takes an exception, re-enters at that same address, and asks the bus for the same word again; the tile answers from the copy it is already holding, so it never observes the repair and spins there, retiring nothing. Recovering that hart requires a power cycle of its tile through the power controller (Section \ref{s:multicore}); \bitfield{ndmreset}, being unimplemented, cannot do it.
	\item \textbf{Debugger software support is not yet available.} No OpenOCD configuration or GDB integration is published for this part. The sequences in Section \ref{ss:debug-worked} are at the DMI level and are what such a configuration would generate.
	\item \textbf{The debug port is not yet part of a signed-off physical implementation.} The five pins exist in the chip-level netlist and in the LQFP-100 package definition, but the current signed-off layout predates them.
\end{itemize}

One further note for anyone comparing this implementation against the RISC-V
debug specification or against a reference simulator. Several behaviours here
differ deliberately from the most widely used reference model, and in every case
this implementation follows the specification's text where the reference does
not: an unrecognised JTAG instruction selects BYPASS rather than leaving the
previous register selected; \bitfield{dmistat} is actually driven, and reports
sticky failure as well as sticky busy; \bitfield{idle} reports the real
requirement instead of zero; an out-of-range hart selection reports
\textit{nonexistent} rather than being clamped to the highest hart; and
\register{dmstatus}.\bitfield{version} reports \texttt{3} for specification 1.0.
A debugger written against the specification will find this part conformant; one
written against a particular simulator's quirks may not.

\fi
